
\begin{frame}
	\frametitle{Outils de formalisation}
	\begin{block}{Idée}
		Dans cette slide, expliquer pourquoi nous faisons appels à des outils
		formels dans nos travaux. %
		Expliquer comment nous proposons de les utiliser pour répondre à nos
		problématiques.
	\end{block}

	\begin{itemize}
		\item Formalisation des objets pour la mise en œuvre de procédures
		      génériques en dimension, notamment grâce aux orbites. %
		\item Formalisation des opérations pour permettre leur analyse. %
		      Étude des impacts d'une opération sur un objet avec la détection des
		      changements topologiques et le suivi des évolutions d'orbites. %
	\end{itemize}

\end{frame}


% \begin{frame}
% 	\frametitle{Cartes généralisées}
% 	\begin{block}{Idée}
% 		Ici, je dois pouvoir reprendre la décomposition présentée dans la
% 		présentation à Marseille. %
% 		Il faudra faire attention à nos pas rentrer dans le détail. %
% 		L'essentiel est de faire comprendre que l'on peut distinguer tout type de
% 		cellule topologique. %
% 	\end{block}
% 	\begin{itemize}
% 		\item illustrer la décomposition d'un objet par ses dimensions
% 		\item introduire les termes brin et orbite
% 	\end{itemize}
% \end{frame}

\begin{frame}{Cartes généralisées (G-cartes)}
	\begin{minipage}[t][.6\textheight][c]{\textwidth}
		\begin{columns}
			\begin{column}{.3\textwidth}
				\begin{overlayarea}{\textwidth}{.6\textheight}
					\begin{figure}
						\centering
						\includegraphics<1>[height=.5\textheight]{decomposition-3D-no-decomp}%
						\includegraphics<2>[height=.5\textheight]{decomposition-3D-3-cell}%
						\includegraphics<3>[height=.5\textheight]{decomposition-3D-2-cell}%
						\includegraphics<4>[height=.5\textheight]{decomposition-3D-1-cell}%
						\includegraphics<5>[height=.5\textheight]{decomposition-3D-0-cell}%
					\end{figure}
				\end{overlayarea}
			\end{column}
			\hfill
			\begin{column}{.4\textwidth}
				\begin{block}<2->{Décomposition d'un objet en :}
					% \begin{overlayarea}{\textwidth}{.3\textheight}
					\begin{itemize}
						\item<2-> Volumes connectés par des $3$-liaisons ({\color{mydarkgreenink} —})%
						\item<3-> Faces connectées par des $2$-liaisons ({\color{myblueink} —})%
						\item<4-> Arêtes connectées par des $1$-liaisons ({\color{red} —})%
						\item<5> Sommets connectés par des $0$-liaisons ({\color{black} —})%
					\end{itemize}
					% \end{overlayarea}
				\end{block}
			\end{column}
		\end{columns}
	\end{minipage}
	\begin{block}{}
		\begin{minipage}[t][.05\textheight][t]{1.0\linewidth}
			Graphe dont les nœuds sont des \emph{brins} et les arcs des \emph{liaisons}.
		\end{minipage}
	\end{block}
\end{frame}

\begin{frame}{Orbites}
	\begin{minipage}[t][.6\textheight][c]{\textwidth}
		\begin{columns}
			\begin{column}{.3\textwidth}
				\begin{overlayarea}{\textwidth}{.6\textheight}
					\begin{figure}
						\includegraphics<1>[height=.5\textheight]{orbit-3D-vertex-123}%
						\includegraphics<2>[height=.5\textheight]{orbit-3D-edge-023}%
						\includegraphics<3>[height=.5\textheight]{orbit-3D-face-013}%
						\includegraphics<4>[height=.5\textheight]{orbit-3D-volume-012}%
						\includegraphics<5>[height=.5\textheight]{orbit-3D-edge-02}%
					\end{figure}
				\end{overlayarea}
			\end{column}
			\hfill
			\begin{column}{.4\textwidth}
				\begin{block}{Orbites~:}
					% \begin{overlayarea}{.5\textwidth}{.45\textheight}
					\begin{itemize}
						\item<1-> Sommet {\color{red}$\orbit{1,2,3}$}{\color{mydarkgreenink}$(a)$}%
						\item<2-> Arête {\color{red}$\orbit{0,2,3}$}{\color{mydarkgreenink}$(a)$}%
						\item<3-> Face {\color{red}$\orbit{0,1,3}$}{\color{mydarkgreenink}$(a)$}%
						\item<4-> Volume {\color{red}$\orbit{0,1,2}$}{\color{mydarkgreenink}$(a)$}%
						\item<5> Arête de volume {\color{red}$\orbit{0,2}$}{\color{mydarkgreenink}$(a)$}
						\item<5> \ldots
					\end{itemize}
					% \end{overlayarea}
				\end{block}
			\end{column}
		\end{columns}
	\end{minipage}
	\begin{block}{}
		\begin{minipage}[t][.05\textheight][t]{1.0\linewidth}
			Sous-graphe de type {\color{red}$\orbit{o}$} incident à un brin
				{\color{mydarkgreenink}(b)}.
		\end{minipage}
	\end{block}
\end{frame}

% \begin{frame}
% 	\frametitle{Règles Jerboa}
% 	\begin{block}{Idée}
% 		Dans cette slide, expliquer que les règles Jerboa rendent possible une
% 		analyse qui n'est pas basée sur du code, mais sur une grammaire formelle. %
% 		L'analyse du code permettant d'arriver au même résultat serait
% 		particulièrement complexe. %
% 		Une grammaire formelle nous permet de repérer plus facilement des entités
% 		topologiques qui sont entre autres créés, supprimés et préservés. %
% 	\end{block}
% 	\begin{itemize}
% 		\item illustrer l'application d'une règle et la transformation qui correspond
% 		\item introduire création, suppression, préservation de nœud, arcs
% 		      implicites et explicites.
% 	\end{itemize}
% \end{frame}

\begin{frame}
	\frametitle{Règles de transformation Jerboa}
	\begin{figure}
		\centering
		\includegraphics<1>[width=.8\textwidth]{triangulation-geometric-view-3D}%
		\includegraphics<2>[height=.8\textheight]{triangulation-application-1-3D}%
		\includegraphics<3>[height=.8\textheight]{triangulation-application-2-left-3D}%
		\includegraphics<4>[height=.8\textheight]{triangulation-application-n0}%
		\includegraphics<5>[height=.8\textheight]{triangulation-application-n2}%
		\includegraphics<6>[height=.8\textheight]{triangulation-application-n1}%
		\includegraphics<7>[height=.8\textheight]{triangulation-application-n0-n1}%
		\includegraphics<8>[height=.8\textheight]{triangulation-application-n1-n2}%
	\end{figure}
\end{frame}


%%% Local Variables:
%%% mode: LaTeX
%%% TeX-master: "../../main"
%%% End:
